\chapter{Performance Evaluation}\label{ch:evaluation}

This chapter will present and discuss the results of several benchmarks which
were carried out to compare the `Vinculum' library presented in \Cref{ch:implementation}
with existing solutions.

\section{Repository Implementation}\label{sec:repo-implementation}

Since Vinculum requires users to provide repository implementations it was
necessary to create one for the benchmarks.
This is done by the `vinculum-benchmark' program which also provides a simple
command line interface to perform the operations of the Lock-Free Deduplication
algorithm.

The repository stores manifest files, chunks and fossils as individual files
in a POSIX~\cite{Posix} conforming file system using the Tokio~\cite{Tokio}
library to perform file system operations in parallel.

The manifest files are directly stored in a directory called `manifests' with their
file name corresponding to the name assigned by the user.
To reduce complexity the list of chunks referenced by them is not stored
in chunks like it is done in Duplicacy~\cite[section `Naming Chunks']{Duplicacy}
but directly stored in the files using a simple binary encoding.

Chunks and fossils are stored in separate directories to allow them to be
enumerated efficiently as required by \Cref{sec:unreferenced-chunks} and
\Cref{sec:abandoned-fossils}.
While some programs like Borg 2~\cite[section `Internals']{Borg2} store them
using a nested directory structure to work around file system limitations they
are stored directly in the directories for simplicity with fossils having the
same name as their associated chunks.

Clients are also stored as files in a separate directory to implement the client
management described in \Cref{sec:adding-clients} and \Cref{sec:removing-clients}.
By storing them as individual files they can be atomically added or removed from
the repository in the same way as chunks and fossils.

To implement the approach described in \Cref{sec:creation} an additional
directory is created containing files which are in the process of being created,
with the temporary files receiving random file names and are opened using the
`O\_EXCL'~\cite{Posix} flag to ensure that clients do not open existing files
when the file name is already in use.

\section{Participating Programs and Environment}

\begin{table}
    \centering
    \begin{tabular}{ c c c c c }
        \toprule
        Program & Borg 1 & Borg 2 & Duplicacy & vinculum-benchmark \\
        \midrule
        Version & 1.4.0  & 2.0.0b14 & 3.2.4 & 0.1.0 \\
        \midrule
        Programming language & Python, C & Python, C & Go & Rust \\
        \bottomrule
    \end{tabular}
    \caption{Programs participating in the benchmarks.}\label{tab:benchmark-programs}
\end{table}

The programs participating in the following benchmarks are all implementations
of the content-addressable storage approach described in \Cref{sec:current-developments}
to ensure a similar internal architecture.
The selection was performed based on the selection used in the paper of Chen, Gilbert
et al.~\cite{Duplicacy} and the ability to be configured with the settings shown in
\Cref{tab:benchmark-settings}.

\begin{table}
    \centering
    \begin{tabular}{ c c }
        \toprule
        Setting & Value \\
        \midrule
        Chunking algorithm & fixed-size \\
        \midrule
        Compression & none \\
        \midrule
        Encryption & none \\
        \bottomrule
    \end{tabular}
    \caption{Settings used in the benchmarks.}\label{tab:benchmark-settings}
\end{table}

These settings were selected to prevent the vinculum-benchmark program developed
to be a simple frontend to Vinculum from having an advantage by not supporting
and therefore having to pay the associated costs of features such as variable-sized
chunking and encryption.
The exact ways in which the settings were applied to the individual programs can be
obtained by reading the benchmark scripts bundled with the source code of Vinculum.

All benchmarks were performed on Linux 6.12.10 running on an AMD Ryzen 5 7600 6-Core
processor with 16 GB of memory.
To reduce disturbances all backups were performed without network access using a
local hard-drive formatted with the ext4 file system and not in use by other programs.

\section{Virtual Machine Images}

\subsection{Setup}

The first benchmark is inspired by one of the benchmarks performed in the paper
of Chen, Gilbert et al.~\cite{Duplicacy}.
The original benchmark consists of creating a VirtualBox Virtual Machine Image
of a machine running CentOS 7 and measuring the creation of backups of it at
the following points~\cite[section `EVALUATION AND ANALYSIS']{Duplicacy}:
\begin{enumerate}
    \item after completing the setup
    \item after installing common developer tools
    \item after a power on immediately followed by a power off
\end{enumerate}

For this benchmark Ubuntu Desktop 24.10 was used and the following APT packages
selected as developer tools:
\begin{itemize}
    \item clang
    \item rust-all
    \item git
    \item vim
    \item cmake
\end{itemize}

Since Vinculum only concerns itself with performing the fossil collection
and deletion steps the benchmark was altered to not measure the creation of
the backups but their removal.

After creating the backups with a fixed chunk size of one MiB they were pruned
in the order of their creation.
All programs required multiple steps to delete a single backup with them being
fossil collection and deletion in the case of Duplicacy and vinculum-benchmark
and deletion and compaction in the case of Borg 1 and 2.
The time required by the individual steps was measured using the elapsed time
between the start of the commands programs and their termination.

\subsection{Results}\label{sec:vm-images-results}

\pgfplotstableread[col sep = comma]{benchmark/results/vm_images.csv}{\vmimagedata}
\begin{figure}[H]
    \centering
    \begin{tikzpicture}
        \pgfplotsset{every axis/.style={
            ybar stacked,
            ymin=0,ymax=170,
            xtick=data,
            xtick style={draw=none},
            symbolic x coords={
                Borg 1,
                Borg 2,
                Duplicacy,
                Vinculum
            },
            bar width=8pt,
            width=.9\textwidth,
            height=.4\textheight
        }}

        \begin{axis}[hide axis,bar shift=-10pt]
            \addplot table [x=tool,y=ubuntu_setup1] {\vmimagedata};
            \addplot table [x=tool,y=ubuntu_setup2] {\vmimagedata};
        \end{axis}

        \begin{axis}[hide axis]
            \addplot table [x=tool,y=ubuntu_tools1] {\vmimagedata};
            \addplot table [x=tool,y=ubuntu_tools2] {\vmimagedata};
        \end{axis}

        \begin{axis}[
            bar shift=10pt,
            xlabel=Backup deletions grouped by program,
            ylabel=Runtime in seconds,
            name=mainplot
        ]
            \addplot table [x=tool,y=ubuntu_onoff1] {\vmimagedata};
            \addplot table [x=tool,y=ubuntu_onoff2] {\vmimagedata};
        \end{axis}

        % subplot
        \begin{axis}[
            hide axis,
            bar shift=-10pt,
            restrict x to domain=1:3,
            at={(mainplot.north east)},
            anchor=north east,
            xshift=-1cm,
            yshift=-1cm,
            ymax=3,
            small
        ]
            \addplot table [x=tool,y=ubuntu_setup1] {\vmimagedata};
            \addplot table [x=tool,y=ubuntu_setup2] {\vmimagedata};
        \end{axis}

        \begin{axis}[
            hide axis,
            restrict x to domain=1:3,
            at={(mainplot.north east)},
            anchor=north east,
            xshift=-1cm,
            yshift=-1cm,
            ymax=3,
            small
        ]
            \addplot table [x=tool,y=ubuntu_tools1] {\vmimagedata};
            \addplot table [x=tool,y=ubuntu_tools2] {\vmimagedata};
        \end{axis}

        \begin{axis}[
            bar shift=10pt,
            restrict x to domain=1:3,
            at={(mainplot.north east)},
            anchor=north east,
            xshift=-1cm,
            yshift=-1cm,
            ymax=3,
            minor y tick num=1,
            small
        ]
            \addplot table [x=tool,y=ubuntu_onoff1] {\vmimagedata};
            \addplot table [x=tool,y=ubuntu_onoff2] {\vmimagedata};
        \end{axis}
    \end{tikzpicture}
    \caption{Time required to delete individual backups.}\label{fig:vm-images}
\end{figure}

\Cref{fig:vm-images} shows the results of performing the benchmark with the time
taken by first step of backup deletion being shown in blue and the second step
shown in red.

All programs exhibit a similar pattern when comparing the total durations of their
deletion operations which can be explained by the amount of data deleted by each
operation.
The first backup was created after completing the setup and shared a large amount
of data with the other two backups, causing only data changed by the installation
of the developer tools to be deleted.
The second backup shared almost all data with the third backup which in turn
caused almost no data to be deleted.
At the end the third backup was the last backup remaining causing all remaining
data to be deleted.

The most noticeable difference is the large amount of time required by Borg 1
in comparison to the other programs with a possible explanation being the way
chunks are stored on disk.
While the other programs store them using individual files Borg 1 combines them
into `segments' which are files behaving like append-only logs~\cite[section `Internals']{Borg1}.
When deleting individual chunks they are only marked as deleted by writing
a special entry into a segment file with the actual deletion being done by
the compaction step.
During this step segments containing a certain number of deleted entries are
read and the remaining entries stored as a new segment using information recorded
during backup creation and deletion~\cite[section `Internals']{Borg1}.
Since this may require additional I/O operations in comparison to deleting individual
files it likely caused a slowdown depending on the amount of data deleted.

The remaining programs featured a comparatively similar performance with Borg 2
being slower than Duplicacy which in turn was slower than Vinculum.
To better show their differences \Cref{fig:vm-images} provides a subplot which
shows them without the results generated by Borg 1.

While the exact causes probably depend on the individual implementations the
different programming languages used by them may be a possible source of
performance differences.
Another reason might be their internal complexity with the support for compression,
encryption and caching provided by Borg 2 and Duplicacy generating some overhead
in comparison to the simple frontend implemented by vinculum-benchmark.

\section{Chunk Scaling}

\subsection{Setup}

The next benchmark is designed to explore how the deletion operations of the
participating programs scale in terms of the number of chunks existing in the
repository.

It consists of first creating a specific number of chunks divided evenly between
four backups with one of them being deleted afterwards.
Since the benchmark in \Cref{sec:backup-scaling} exhibited a measurable
slowdown caused by excessive I/O operations related to caching the cache of Borg 2
and Duplicacy is located on a tmpfs for the reasons explained there.

The time required by each deletion operation was measured by using the elapsed
time between the start of the individual commands and their termination and
adding the runtimes of the individual steps (such as the fossil collection and
deletions steps) together.

\subsection{Results}

There exists again a similarity between the behaviors of the participating
programs with their runtime increasing almost linearly in terms of the number
of chunks, particularly clearly visible by the increased number of measurements
depicted in \Cref{fig:chunks-large}.

While Borg 1 does not exhibit a significantly increased runtime when using small
chunk sizes its performance degrades when increasing them as shown in \Cref{fig:chunk-scaling}.
This is again likely caused by its segment files which provide an advantage when
the data contained in them is small by reducing the amount of file system operations
such as deletions to be performed.
Should their size increase the overhead created by the compaction algorithm described
in \Cref{sec:vm-images-results} will also increase and seems to surpass the advantage of
having to manage a smaller number of files, turning it into a disadvantage.

While the chunk size doesn't seem to affect the rest of the programs for the most
part its shown that Duplicacy and Vinculum feature an almost identical increase
in runtime under the circumstances.
While there is a slight difference causing Vinculum to be faster than Duplicacy
its constant nature suggests it being caused by something independent of the
implementations of the fossil collection and deletion steps such as internal
complexity or time required for startup.

\begin{figure}[H]
    \centering
    \foreach \f in {64, 4096}{
        \pgfplotstableread[col sep = comma]{benchmark/results/chunks\f .csv}{\chunkdata}

        \begin{subfigure}{\textwidth}
            \centering
            \begin{tikzpicture}
                \pgfplotsset{
                    width=.9\textwidth
                }
                \begin{axis}[
                    xlabel=Number of chunks,
                    ylabel=Runtime in seconds,
                    width=0.75\textwidth,
                    legend pos=outer north east,
                    scaled x ticks=base 10:-3
                ]
                    \addplot table [x=chunks, y expr=\thisrow{Borg 1 deletion}+\thisrow{Borg 1 compaction}] {\chunkdata};
                    \addlegendentry{Borg 1}
                    \addplot table [x=chunks, y expr=\thisrow{Borg 2 deletion}+\thisrow{Borg 2 compaction}] {\chunkdata};
                    \addlegendentry{Borg 2}
                    \addplot table [x=chunks, y expr=\thisrow{Duplicacy collection}+\thisrow{Duplicacy deletion}] {\chunkdata};
                    \addlegendentry{Duplicacy}
                    \addplot table [x=chunks, y expr=\thisrow{Vinculum collection}+\thisrow{Vinculum deletion}] {\chunkdata};
                    \addlegendentry{Vinculum}
                \end{axis}
            \end{tikzpicture}
            \caption{chunk size = \f{} bytes.}\label{fig:chunks\f}
        \end{subfigure}
    }
    \caption{Time required to delete a single backup with different chunk sizes.}
\end{figure}

\begin{figure}[H]\ContinuedFloat
    \foreach \f in {32768, 524288}{
        \pgfplotstableread[col sep = comma]{benchmark/results/chunks\f .csv}{\chunkdata}

        \begin{subfigure}{\textwidth}
            \centering
            \begin{tikzpicture}
                \pgfplotsset{
                    width=.9\textwidth
                }
                \begin{axis}[
                    xlabel=Number of chunks,
                    ylabel=Runtime in seconds,
                    width=0.75\textwidth,
                    legend pos=outer north east,
                    scaled x ticks=base 10:-3
                ]
                    \addplot table [x=chunks, y expr=\thisrow{Borg 1 deletion}+\thisrow{Borg 1 compaction}] {\chunkdata};
                    \addlegendentry{Borg 1}
                    \addplot table [x=chunks, y expr=\thisrow{Borg 2 deletion}+\thisrow{Borg 2 compaction}] {\chunkdata};
                    \addlegendentry{Borg 2}
                    \addplot table [x=chunks, y expr=\thisrow{Duplicacy collection}+\thisrow{Duplicacy deletion}] {\chunkdata};
                    \addlegendentry{Duplicacy}
                    \addplot table [x=chunks, y expr=\thisrow{Vinculum collection}+\thisrow{Vinculum deletion}] {\chunkdata};
                    \addlegendentry{Vinculum}
                \end{axis}
            \end{tikzpicture}
            \caption{chunk size = \f{} bytes.}\label{fig:chunks\f}
        \end{subfigure}
    }
    \caption{Time required to delete a single backup with different chunk sizes.}\label{fig:chunk-scaling}
\end{figure}

\begin{figure}
    \centering
    \pgfplotstableread[col sep = comma]{benchmark/results/chunks4096large.csv}{\chunkdata}
    \begin{tikzpicture}
        \pgfplotsset{
            width=.9\textwidth
        }
        \begin{axis}[
            xlabel=Number of chunks,
            ylabel=Runtime in seconds,
            width=0.75\textwidth,
            legend pos=outer north east,
            scaled x ticks=base 10:-3
        ]
            \addplot table [x=chunks, y expr=\thisrow{Borg 1 deletion}+\thisrow{Borg 1 compaction}] {\chunkdata};
            \addlegendentry{Borg 1}
            \addplot table [x=chunks, y expr=\thisrow{Borg 2 deletion}+\thisrow{Borg 2 compaction}] {\chunkdata};
            \addlegendentry{Borg 2}
            \addplot table [x=chunks, y expr=\thisrow{Duplicacy collection}+\thisrow{Duplicacy deletion}] {\chunkdata};
            \addlegendentry{Duplicacy}
            \addplot table [x=chunks, y expr=\thisrow{Vinculum collection}+\thisrow{Vinculum deletion}] {\chunkdata};
            \addlegendentry{Vinculum}
        \end{axis}
    \end{tikzpicture}
    \caption{Further measurements of the time required to delete a single backup with chunk size = 4096 bytes.}\label{fig:chunks-large}
\end{figure}

\section{Backups Scaling}\label{sec:backup-scaling}

\subsection{Setup}\label{sec:backup-scaling-setup}

While the previous benchmark explored how the deletion operations of the
participating programs scaled in terms of the number of chunks existing in the
repository the next benchmark will do so in terms of the number of backups.

It consists of creating a fixed number of chunks with a size of 4096 bytes which
are evenly divided between a specific number of backups with one of them being
deleted afterwards.
Since especially Duplicacy exhibited a measurable slowdown caused by excessive I/O
operations related to caching the cache of Borg 2 and Duplicacy is located on a
tmpfs to prevent Vinculum from gaining an advantage by not implementing caching.

The time required by each deletion operation was measured by using the elapsed
time between the start of the individual commands and their termination and
adding the runtimes of the individual steps (such as the fossil collection and
deletions steps) together.

\subsection{Results}

\begin{figure}
    \foreach \f in {0, 1}{
        \pgfplotstableread[col sep = comma]{benchmark/results/manifests\f const.csv}{\manifestdata}

        \begin{subfigure}{\textwidth}
            \centering
            \begin{tikzpicture}
                \pgfplotsset{
                    width=.9\textwidth
                }
                \begin{axis}[
                    xlabel=Number of backups,
                    ylabel=Runtime in seconds,
                    width=0.75\textwidth,
                    legend pos=outer north east,
                    name=mainplot
                ]
                    \addplot table [x=manifests, y expr=\thisrow{Borg 1 deletion}+\thisrow{Borg 1 compaction}] {\manifestdata};
                    \addlegendentry{Borg 1}
                    \addplot table [x=manifests, y expr=\thisrow{Borg 2 deletion}+\thisrow{Borg 2 compaction}] {\manifestdata};
                    \addlegendentry{Borg 2}
                    \addplot table [x=manifests, y expr=\thisrow{Duplicacy collection}+\thisrow{Duplicacy deletion}] {\manifestdata};
                    \addlegendentry{Duplicacy}
                    \addplot table [x=manifests, y expr=\thisrow{Vinculum collection}+\thisrow{Vinculum deletion}] {\manifestdata};
                    \addlegendentry{Vinculum}
                \end{axis}

                %subplot
                \begin{axis}[
                    cycle list shift=2,
                    at={(mainplot.north west)},
                    anchor=north west,
                    xshift=1cm,
                    yshift=-1cm,
                    scaled x ticks=base 10:-2,
                    footnotesize
                ]
                    \addplot table [x=manifests, y expr=\thisrow{Duplicacy collection}+\thisrow{Duplicacy deletion}] {\manifestdata};
                    \addplot table [x=manifests, y expr=\thisrow{Vinculum collection}+\thisrow{Vinculum deletion}] {\manifestdata};
                \end{axis}
            \end{tikzpicture}
            \caption{total chunks = \f{}.}\label{fig:backups\f}
        \end{subfigure}
    }
    \caption{Time required to delete a single backup with different amounts of chunks.}
\end{figure}

\begin{figure}\ContinuedFloat
    \foreach \f in {1000}{
        \pgfplotstableread[col sep = comma]{benchmark/results/manifests\f const.csv}{\manifestdata}

        \begin{subfigure}{\textwidth}
            \centering
            \begin{tikzpicture}
                \pgfplotsset{
                    width=.9\textwidth
                }
                \begin{axis}[
                    xlabel=Number of backups,
                    ylabel=Runtime in seconds,
                    width=0.75\textwidth,
                    legend pos=outer north east,
                    name=mainplot
                ]
                    \addplot table [x=manifests, y expr=\thisrow{Borg 1 deletion}+\thisrow{Borg 1 compaction}] {\manifestdata};
                    \addlegendentry{Borg 1}
                    \addplot table [x=manifests, y expr=\thisrow{Borg 2 deletion}+\thisrow{Borg 2 compaction}] {\manifestdata};
                    \addlegendentry{Borg 2}
                    \addplot table [x=manifests, y expr=\thisrow{Duplicacy collection}+\thisrow{Duplicacy deletion}] {\manifestdata};
                    \addlegendentry{Duplicacy}
                    \addplot table [x=manifests, y expr=\thisrow{Vinculum collection}+\thisrow{Vinculum deletion}] {\manifestdata};
                    \addlegendentry{Vinculum}
                \end{axis}

                %subplot
                \begin{axis}[
                    cycle list shift=2,
                    at={(mainplot.north west)},
                    anchor=north west,
                    xshift=1cm,
                    yshift=-1cm,
                    scaled x ticks=base 10:-2,
                    footnotesize
                ]
                    \addplot table [x=manifests, y expr=\thisrow{Duplicacy collection}+\thisrow{Duplicacy deletion}] {\manifestdata};
                    \addplot table [x=manifests, y expr=\thisrow{Vinculum collection}+\thisrow{Vinculum deletion}] {\manifestdata};
                \end{axis}
            \end{tikzpicture}
            \caption{total chunks = \f{}.}\label{fig:backups\f}
        \end{subfigure}
    }
    \caption{Time required to delete a single backup with different amounts of chunks.}\label{fig:backups-scaling}
\end{figure}

In contrast to my expectations the results shown in \Cref{fig:backups-scaling}
do not seem to be affected by the total number of chunks existing in the repository
with a possible explanation being the number of chunks being too low to have a
visible effect.

What can be deduced is that while the runtime of Borg 1 and 2 seems to exhibit
quadratic growth the runtime of Duplicacy and Vinculum seems to be constant.
When looking at the magnified data presented in the subplots of \Cref{fig:backups-scaling}
it becomes apparent that their runtime also increases with the number of backups,
albeit much slower and apparently linearly.

The difference between the slopes of the graphs of Duplicacy and Vinculum
may be explained by the cache activity performed by Duplicacy as described
in \Cref{sec:backup-scaling-setup}.
Since Duplicacy needs to access some chunks to download all manifest
files~\cite[section `Naming Chunks']{Duplicacy} and stores them in a cache on the
local file system~\cite[topic `Cache folder is is extremely big!']{Duplicacy-Forum}
there might be additional overhead compared to Vinculum which does not perform
manifest file chunking or caching.
